\section{David Tong QFT PS3}
\subsection{Weyl and Dirac Representations of the Clifford Algebra, Q1}

The Weyl representation of the Clifford algebra is given by
\begin{equation}
    \gamma^0 = \begin{pmatrix} 0 & \mathbf{1} \\ \mathbf{1} & 0 \end{pmatrix}, \qquad
    \gamma^i = \begin{pmatrix} 0 & \sigma^i \\ -\sigma^i & 0 \end{pmatrix}
\end{equation}
Show that these indeed satisfy $\{\gamma^\mu, \gamma^\nu\} = 2\eta^{\mu\nu}$, where $\mathbf{1}$ comes with an implicit $4\times 4$ unit matrix. Find a unitary matrix $U$ such that $(\gamma')^\mu = U\gamma^\mu U^\dagger$, where $(\gamma')^\mu$ form the Dirac representation of the Clifford algebra
\begin{equation}
    (\gamma')^0 = \begin{pmatrix} \mathbf{1} & 0 \\ 0 & -\mathbf{1} \end{pmatrix}, \qquad
    (\gamma')^i = \begin{pmatrix} 0 & \sigma^i \\ -\sigma^i & 0 \end{pmatrix}.
\end{equation}
\begin{solution}
    Write $\sigma^\mu = (\mathbf{1},\mathbf{\sigma})$, and $\bar{\sigma}^\mu = (\mathbf{1},-\mathbf{\sigma})$. Th anticommutator is given by 
    \begin{equation}
    \begin{aligned}
        \{\gamma^\mu,\gamma^\nu\} &= \begin{pmatrix}
            0&\sigma^\mu\\
            \bar{\sigma}^\mu &0
        \end{pmatrix}\begin{pmatrix}
            0&\sigma^\nu\\
            \bar{\sigma}^\nu &0
        \end{pmatrix}+\begin{pmatrix}
            0&\sigma^\nu\\
            \bar{\sigma}^\nu &0
        \end{pmatrix}\begin{pmatrix}
            0&\sigma^\mu\\
            \bar{\sigma}^\mu &0
        \end{pmatrix}\\
        &=\begin{pmatrix}
            \sigma^\mu\bar{\sigma}^\nu + \sigma^\nu\bar{\sigma}^\mu&0\\
            0&\bar{\sigma}^\mu\sigma^\nu + \bar{\sigma}^\nu\sigma^\mu
        \end{pmatrix}
        \\&=2\eta^{\mu\nu}\mathbf{1}.
    \end{aligned}
    \end{equation}
    $U$ diagnolise $\gamma^0$, and left the rest unchanged. Therefore, we find
    \begin{equation}
        U = \frac{1}{\sqrt{2}}\begin{pmatrix}
            \mathbf{1}&\mathbf{1}\\
            \mathbf{1}&-\mathbf{1}
        \end{pmatrix}.
    \end{equation}
\end{solution}

\subsection{The Lorentz Algebra from Gamma Matrices, Q2}

Show that if $\{\gamma^\mu, \gamma^\nu\} = 2\eta^{\mu\nu}$, then
\begin{equation}
    [\gamma^\kappa\gamma^\lambda, \gamma^\mu\gamma^\nu] = 2\eta^{\lambda\mu}\gamma^\kappa\gamma^\nu - 2\eta^{\kappa\mu}\gamma^\lambda\gamma^\nu + 2\eta^{\lambda\nu}\gamma^\mu\gamma^\kappa - 2\eta^{\kappa\nu}\gamma^\mu\gamma^\lambda.
\end{equation}
Show further that $S^{\mu\nu} \equiv \frac{1}{4}[\gamma^\mu, \gamma^\nu] = \frac{1}{2}(\gamma^\mu\gamma^\nu - \eta^{\mu\nu})$. Use this to confirm that the matrices $S^{\mu\nu}$ form a representation of the Lie algebra of the Lorentz group.

\begin{solution}
    Note that
    \begin{equation}
        [\gamma^\kappa\gamma^\lambda,\gamma^\mu\gamma^\nu] = \gamma^\kappa\gamma^\lambda\gamma^\mu\gamma^\nu - \gamma^\mu\gamma^\nu\gamma^\kappa\gamma^\lambda.
    \end{equation}
    Then, 
    Hence, we find
    \begin{equation}
        \begin{aligned}
            \gamma^\kappa\gamma^\lambda\gamma^\mu\gamma^\nu &= 2\eta^{\lambda\mu}\gamma^\kappa\gamma^\nu -   \gamma^\kappa\gamma^\mu\gamma^\lambda\gamma^\nu \\
            &=2\eta^{\lambda\mu}\gamma^\kappa\gamma^\nu -2\eta^{\lambda\nu}\gamma^\kappa\gamma^\mu + \gamma^\kappa\gamma^\mu\gamma^\nu\gamma^\lambda\\
            &=2\eta^{\lambda\mu}\gamma^\kappa\gamma^\nu -2\eta^{\lambda\nu}\gamma^\kappa\gamma^\mu +2\eta^{\kappa\mu}\gamma^\nu\gamma^\lambda -2\eta^{\lambda\nu}\gamma^\mu\gamma^\lambda + \gamma^\mu\gamma^\nu\gamma^\kappa\gamma^\lambda.
            \end{aligned}
    \end{equation}
    Hence, we find
    \begin{equation}
        [\gamma^\kappa\gamma^\lambda, \gamma^\mu\gamma^\nu] = 2\eta^{\lambda\mu}\gamma^\kappa\gamma^\nu - 2\eta^{\kappa\mu}\gamma^\lambda\gamma^\nu + 2\eta^{\lambda\nu}\gamma^\mu\gamma^\kappa - 2\eta^{\kappa\nu}\gamma^\mu\gamma^\lambda.
    \end{equation}

    Similarly,
    \begin{equation}
        \begin{aligned}
            S^{\mu\nu} &= \frac{1}{4}(\gamma^\mu\gamma^\nu -\gamma^\nu\gamma ^\mu  )\\
            &=\frac{1}{4}(2\gamma^\mu\gamma^\nu - 2\eta^{\mu\nu})\\
            &=\frac{1}{2}(\gamma^\mu\gamma^\nu - \eta^{\mu\nu}).
        \end{aligned}
    \end{equation}
    The Lie algebra is given by
        \begin{equation}
    \begin{aligned}
        [S^{\kappa\lambda},S^{\mu\nu}] &= \frac{1}{4}(\gamma^\kappa\gamma^\lambda\gamma^\mu\gamma^\nu -\eta^{\kappa\lambda}\gamma^\mu\gamma^\nu-\eta^{\mu\nu}\gamma^\kappa\gamma^\lambda - \gamma^\mu\gamma^\nu\gamma^\kappa\gamma^\lambda+\eta^{\kappa\lambda}\gamma^\mu\gamma^\nu+\eta^{\mu\nu}\gamma^\kappa\gamma^\lambda)\\
        &=\frac{1}{4}[\gamma^\kappa\gamma^\lambda,\gamma^\mu\gamma^\nu]\\
        &=\frac{1}{4}\left(2\eta^{\lambda\mu}\gamma^\kappa\gamma^\nu - 2\eta^{\kappa\mu}\gamma^\lambda\gamma^\nu + 2\eta^{\lambda\nu}\gamma^\mu\gamma^\kappa - 2\eta^{\kappa\nu}\gamma^\mu\gamma^\lambda\right)\\
        &=\frac{1}{2}\left(\eta^{\lambda\mu}\gamma^\kappa\gamma^\nu - \eta^{\kappa\mu}\gamma^\lambda\gamma^\nu + \eta^{\lambda\nu}\gamma^\mu\gamma^\kappa - \eta^{\kappa\nu}\gamma^\mu\gamma^\lambda\right)\\
        &=\eta^{\lambda\mu}\left(S^{\kappa\nu}+\tfrac{1}{2}\eta^{\kappa\nu}\right) - \eta^{\kappa\mu}\left(S^{\lambda\nu}+\tfrac{1}{2}\eta^{\lambda\nu}\right) + \eta^{\lambda\nu}\left(S^{\mu\kappa}+\tfrac{1}{2}\eta^{\mu\kappa}\right) - \eta^{\kappa\nu}\left(S^{\mu\lambda}+\tfrac{1}{2}\eta^{\mu\lambda}\right)\\
        &= \eta^{\lambda\mu}S^{\kappa\nu} - \eta^{\kappa\mu}S^{\lambda\nu} + \eta^{\lambda\nu}S^{\mu\kappa} - \eta^{\kappa\nu}S^{\mu\lambda}.
    \end{aligned}
    \end{equation}
\end{solution}

\subsection{Trace Technology and Gamma Matrix Identities, Q3}

Using just the algebra $\{\gamma^\mu, \gamma^\nu\} = 2\eta^{\mu\nu}$ (i.e.\ without resorting to a particular representation), and defining $\gamma^5 = -i\gamma^0\gamma^1\gamma^2\gamma^3$, $\slashed{p} = p_\mu\gamma^\mu$ and $S^{\mu\nu} = \frac{1}{4}[\gamma^\mu, \gamma^\nu]$, prove the following results: (\textit{Some useful tricks include the cyclicity of the trace, and inserting $(\gamma^5)^2 = 1$ into a trace}).
\begin{enumerate}[label=\roman*.]
    \item $\mathrm{Tr}\,\gamma^\mu = 0$
    \item $\mathrm{Tr}(\gamma^\mu\gamma^\nu) = 4\eta^{\mu\nu}$
    \item $\mathrm{Tr}(\gamma^\mu\gamma^\nu\gamma^\rho) = 0$
    \item $(\gamma^5)^2 = 1$
    \item $\mathrm{Tr}\,\gamma^5 = 0$
    \item $\slashed{p}\slashed{q} = 2p\cdot q - \slashed{q}\slashed{p} = p\cdot q + 2S^{\mu\nu}p_\mu q_\nu$
    \item $\mathrm{Tr}(\slashed{p}\slashed{q}) = 4p\cdot q$
    \item $\mathrm{Tr}(\slashed{p}_1\cdots\slashed{p}_n) = 0$ if $n$ is odd
    \item $\mathrm{Tr}(\slashed{p}_1\slashed{p}_2\slashed{p}_3\slashed{p}_4) = 4\left[(p_1\cdot p_2)(p_3\cdot p_4) + (p_1\cdot p_4)(p_2\cdot p_3) - (p_1\cdot p_3)(p_2\cdot p_4)\right]$
    \item $\mathrm{Tr}(\gamma^5\slashed{p}_1\slashed{p}_2) = 0$
    \item $\gamma_\mu\slashed{p}\gamma^\mu = -2\slashed{p}$
    \item $\gamma_\mu\slashed{p}_1\slashed{p}_2\gamma^\mu = 4p_1\cdot p_2$
    \item $\gamma_\mu\slashed{p}_1\slashed{p}_2\slashed{p}_3\gamma^\mu = -2\slashed{p}_3\slashed{p}_2\slashed{p}_1$
    \item $\mathrm{Tr}(\gamma^5\slashed{p}_1\slashed{p}_2\slashed{p}_3\slashed{p}_4) = 4i\,\epsilon_{\mu\nu\rho\sigma}\,p_1^\mu p_2^\nu p_3^\rho p_4^\sigma$
\end{enumerate}

\begin{solution}
Let us consider spinors in $d$-dimensions.
    \begin{enumerate}[label=\roman*.]
    \item Trivial.
    \item 
    \begin{equation}
            \Tr{\gamma^\mu\gamma^\nu} = 2\eta^{\mu\nu}\Tr{\mathbf{1}} -\Tr{\gamma^\nu\gamma^\mu}\implies\Tr{\gamma^\mu\gamma^\nu} = d\eta^{\mu\nu}.
        \end{equation}
    \item 
    \begin{equation}
        \Tr{\gamma^\mu\gamma^\nu\gamma^\rho} = -\Tr{\gamma^\nu\gamma^\mu\gamma^\rho} =\Tr{\gamma^\nu\gamma^\rho\gamma^\mu} = -\Tr{\gamma^\rho\gamma^\nu\gamma^\mu} =0.
    \end{equation}
    \item 
    \begin{equation}
        (\gamma^5)^2 = -\gamma^0\gamma^1\gamma^2\gamma^3\gamma^0\gamma^1\gamma^2\gamma^3=-\gamma^1\gamma^2\gamma^3\gamma^1\gamma^2\gamma^3 =-\gamma^2\gamma^3\gamma^2\gamma^3 =\mathbf{1}.
    \end{equation}
    \item Will turn out to be trivial using the general trace of for $\gamma$-matrices.
    \item 
    \begin{equation}
        \slashed{p}\slashed{q} = p_\mu q_\nu(2\eta^{\mu\nu}-\gamma^\nu\gamma^\mu) = 2p\cdot q - \slashed{q}\slashed{p} = p_\mu q_\nu(\eta^{\mu\nu}+2S^{\mu\nu}) = p\cdot q + 2S^{\mu\nu}p_\mu q_\nu.
    \end{equation}
    \item 
    \begin{equation}
        \Tr{\slashed{p}\slashed{q}} = d\,p\cdot q.
    \end{equation}
    \item Follows from iii by induction.
    \item 
    \begin{equation}
        \begin{aligned}
            &\Tr{\gamma^\mu\gamma^\nu\gamma^\rho\gamma^\lambda} = 2(\eta^{\mu\nu}\Tr{\gamma^\rho\gamma^\lambda}-\eta^{\mu\rho}\Tr{\gamma^\nu\gamma^\lambda}+\eta^{\mu\lambda}\Tr{\gamma^\nu\gamma^\rho}) -\Tr{\gamma^\nu\gamma^\rho\gamma^\lambda\gamma^\mu}\\
            \implies& \Tr{\gamma^\mu\gamma^\nu\gamma^\rho\gamma^\lambda} = d(\eta^{\mu\nu}\eta^{\rho\lambda}-\eta^{\mu\rho}\eta^{\nu\lambda}+\eta^{\mu\lambda}\eta^{\nu\rho})\\
            \implies &\Tr{\slashed{p}_1\slashed{p}_2\slashed{p}_3\slashed{p}_4} = d\left[(p_1\cdot p_2)(p_3\cdot p_4) + (p_1\cdot p_4)(p_2\cdot p_3) - (p_1\cdot p_3)(p_2\cdot p_4)\right].
        \end{aligned}
    \end{equation}
    \item 
    A cleaner argument: by anticommutation of $\gamma^5$ with all $\gamma^\mu$, the trace is totally antisymmetric in $(\mu, \nu)$. The only available rank-2 tensor in 4d is $\eta^{\mu\nu}$, which is symmetric. Hence the trace must vanish:
    \begin{equation}
        \Tr{\gamma^5\slashed{p}_1\slashed{p}_2} = 0.
    \end{equation}
    \item Insert $\eta_{\mu\nu}\gamma^\mu\gamma^\nu = d\,\mathbf{1}$ and use $\{\gamma^\mu, \gamma^\nu\} = 2\eta^{\mu\nu}$:
    \begin{equation}
        \gamma_\mu\slashed{p}\gamma^\mu = p_\nu\gamma_\mu\gamma^\nu\gamma^\mu = p_\nu(2\eta^\nu_\mu - \gamma^\nu\gamma_\mu)\gamma^\mu = 2\slashed{p} - d\slashed{p} = (2-d)\slashed{p} = -2\slashed{p}.
    \end{equation}
    
    \item Apply the previous result twice. Using $\gamma_\mu\gamma^\nu\gamma^\rho\gamma^\mu = 4\eta^{\nu\rho}$ (proved similarly):
    \begin{equation}
        \gamma_\mu\slashed{p}_1\slashed{p}_2\gamma^\mu = p_{1\nu}p_{2\rho}\gamma_\mu\gamma^\nu\gamma^\rho\gamma^\mu = 4 p_{1\nu}p_{2\rho}\eta^{\nu\rho} = 4 p_1\cdot p_2.
    \end{equation}
    
    \item Similarly $\gamma_\mu\gamma^\nu\gamma^\rho\gamma^\sigma\gamma^\mu = -2\gamma^\sigma\gamma^\rho\gamma^\nu$, hence
    \begin{equation}
        \gamma_\mu\slashed{p}_1\slashed{p}_2\slashed{p}_3\gamma^\mu = -2\slashed{p}_3\slashed{p}_2\slashed{p}_1.
    \end{equation}
    
    \item By the same antisymmetry argument as in (x), the trace $\Tr{\gamma^5\slashed{p}_1\slashed{p}_2\slashed{p}_3\slashed{p}_4}$ must be totally antisymmetric in the four momentum indices. The only such rank-4 tensor in 4d is $\epsilon^{\mu\nu\rho\sigma}$, so
    \begin{equation}
        \Tr{\gamma^5\gamma^\mu\gamma^\nu\gamma^\rho\gamma^\sigma} = c\,\epsilon^{\mu\nu\rho\sigma}.
    \end{equation}
    The constant $c$ is fixed by setting $(\mu,\nu,\rho,\sigma) = (0,1,2,3)$:
    \begin{equation}
        \Tr{\gamma^5\gamma^0\gamma^1\gamma^2\gamma^3} = i\Tr{\gamma^5\gamma^5} = 4i,
    \end{equation}
    using $\gamma^5 = -i\gamma^0\gamma^1\gamma^2\gamma^3$ and $(\gamma^5)^2 = \mathbf{1}$. With $\epsilon^{0123} = -\epsilon_{0123} = -1$, we find $c = -4i$, giving
    \begin{equation}
        \Tr{\gamma^5\slashed{p}_1\slashed{p}_2\slashed{p}_3\slashed{p}_4} = -4i\,\epsilon_{\mu\nu\rho\sigma}\,p_1^\mu p_2^\nu p_3^\rho p_4^\sigma.
    \end{equation}
    \end{enumerate}
\end{solution}

\subsection{Plane-Wave Spinor Solutions and Orthogonality, Q4}

The plane-wave solutions to the Dirac equation are
\begin{equation}
    u^s(\mathbf{p}) = \begin{pmatrix}\sqrt{p\cdot\sigma}\,\xi^s \\ \sqrt{p\cdot\bar\sigma}\,\xi^s\end{pmatrix} \quad \text{and} \quad v^s(\mathbf{p}) = \begin{pmatrix}\sqrt{p\cdot\sigma}\,\xi^s \\ -\sqrt{p\cdot\bar\sigma}\,\xi^s\end{pmatrix}
\end{equation}
where $\sigma^\mu = (1, \mathbf{\sigma}^i)$ and $\bar\sigma^\mu = (1, -\mathbf{\sigma}^i)$ and $\xi^s$, with $s=1,2$, is a basis of orthonormal two-component spinors, satisfying $(\xi^r)^\dagger\cdot\xi^s = \delta^{rs}$. Show that
\begin{align}
    u^r(\mathbf{p})^\dagger\cdot u^s(\mathbf{p}) &= 2p_0\delta^{rs}\\
    \bar{u}^r(\mathbf{p})\cdot u^s(\mathbf{p}) &= 2m\delta^{rs}
\end{align}
and similarly,
\begin{align}
    v^r(\mathbf{p})^\dagger\cdot v^s(\mathbf{p}) &= 2p_0\delta^{rs}\\
    \bar{v}^r(\mathbf{p})\cdot v^s(\mathbf{p}) &= -2m\delta^{rs}
\end{align}
Show also that the orthonality condition between $u$ and $v$ is
\begin{equation}
    \bar{u}^s(\mathbf{p})\cdot v^r(\mathbf{p}) = 0
\end{equation}
while taking the inner product using $\dagger$ requires an extra minus sign
\begin{equation}
    u^r(\mathbf{p})^\dagger\cdot v^s(-\mathbf{p}) = 0.
\end{equation}

\begin{solution}
    \begin{equation}
    \begin{aligned}
     u^r(\mathbf{p})^\dagger\cdot u^s(\mathbf{p}) &= \begin{pmatrix}
         \sqrt{p\cdot \sigma^\dagger}\xi^{\dagger r}&\sqrt{p\cdot\bar{\sigma}^\dagger}\xi^{\dagger r}
     \end{pmatrix}
     \begin{pmatrix}
         \sqrt{p\cdot \sigma}\xi^s\\\sqrt{p\cdot\bar{\sigma}}\xi^s
     \end{pmatrix}\\
     &=\delta^{rs}\left[\sqrt{(p\cdot\sigma)(p\cdot\sigma^\dagger)} + \sqrt{(p\cdot\bar\sigma)(p\cdot\bar\sigma^\dagger)}\right]\\
     &=2\delta^{rs}\sqrt{(p\cdot\sigma)(p\cdot\sigma^\dagger)},
    \end{aligned}
    \end{equation}    
    Since the Pauli matrices are Hermitian, we find
    \begin{equation}
        u^r(\mathbf{p})^\dagger\cdot u^s(\mathbf{p}) = 2\delta^{rs}(p\cdot\sigma +p\cdot\bar\sigma ) = 2p_0\delta^{rs}.
    \end{equation}
    Similarly, we find
    \begin{equation}
        \begin{aligned}
        \bar u^r(\mathbf{p})\cdot u^s(\mathbf{p}) &= \begin{pmatrix}
         \sqrt{p\cdot\bar{\sigma}^\dagger}\xi^{\dagger r}&\sqrt{p\cdot \sigma^\dagger}\xi^{\dagger r}
     \end{pmatrix}
     \begin{pmatrix}
         \sqrt{p\cdot \sigma}\xi^s\\\sqrt{p\cdot\bar{\sigma}}\xi^s
     \end{pmatrix}\\
     &=\delta^{rs}\left[\sqrt{(p\cdot\sigma)(p\cdot\bar\sigma^\dagger)} + \sqrt{(p\cdot\bar\sigma)(p\cdot\sigma^\dagger)}\right]\\
     &=2\delta^{rs}\sqrt{(p\cdot\sigma)(p\cdot\bar\sigma)}\\
     &=2m\delta^{rs},
    \end{aligned}
    \end{equation}
    where we have expanded $p\cdot\sigma$ and $p\cdot\bar\sigma$ and evaluated them explicitly.

    Similarly, one can easily obtain
    \begin{align}
    v^r(\mathbf{p})^\dagger\cdot v^s(\mathbf{p}) &= 2p_0\delta^{rs}\\
    \bar{v}^r(\mathbf{p})\cdot v^s(\mathbf{p}) &= -2m\delta^{rs}.
    \end{align}

    Now, we demonstrate the orthogonality condition,
    \begin{equation}
    \begin{aligned}
            \bar{u}^s(\mathbf{p})\cdot v^r(\mathbf{p})&=\begin{pmatrix}
         \sqrt{p\cdot\bar{\sigma}^\dagger}\xi^{\dagger r}&\sqrt{p\cdot \sigma^\dagger}\xi^{\dagger r}
     \end{pmatrix}
     \begin{pmatrix}
         \sqrt{p\cdot \sigma}\xi^s\\-\sqrt{p\cdot\bar{\sigma}}\xi^s
     \end{pmatrix}=0,
    \end{aligned}
    \end{equation}
    by Hermiticity of Pauli matrices. Similarly,
    \begin{equation}
        u^r(\mathbf{p})^\dagger\cdot v^s(-\mathbf{p}) = \begin{pmatrix}
         \sqrt{p\cdot \sigma^\dagger}\xi^{\dagger r}&\sqrt{p\cdot\bar{\sigma}^\dagger}\xi^{\dagger r}
     \end{pmatrix}
     \begin{pmatrix}
         \sqrt{ \tilde p\cdot \sigma}\xi^s\\-\sqrt{\tilde p\cdot\bar{\sigma}}\xi^s
     \end{pmatrix},
    \end{equation}
    where $\tilde p^\mu = (p_0,-\mathbf{p})$. Now, note that $\tilde p\cdot \sigma = p\cdot \bar\sigma$, we see that the above inner product vanishes.
\end{solution}

\subsection{Spin Sums and Completeness Relation, Q5}

Using the same notation as Question 4 show that
\begin{align}
    \sum_{s=1}^{2} u^s(\mathbf{p})\bar{u}^s(\mathbf{p}) &= \slashed{p} + m \label{eq:uouter}\\
    \sum_{s=1}^{2} v^s(\mathbf{p})\bar{v}^s(\mathbf{p}) &= \slashed{p} - m \label{eq:vouter}
\end{align}
where, rather than being contracted, the two spinors on the left-hand side are placed back to back to form a $4\times 4$ matrix.

\begin{solution}
    \begin{equation}
        \begin{aligned}
            \sum_{s=1}^{2} u^s(\mathbf{p})\bar{u}^s(\mathbf{p}) &= \sum_{s=1}^2
            \begin{pmatrix}
         \sqrt{p\cdot \sigma}\xi^s\\\sqrt{p\cdot\bar{\sigma}}\xi^s
     \end{pmatrix}\begin{pmatrix}
         \sqrt{p\cdot\bar{\sigma}^\dagger}\xi^{\dagger s}&\sqrt{p\cdot \sigma^\dagger}\xi^{\dagger s}
     \end{pmatrix}
     \\
     &=\begin{pmatrix}
         m&p\cdot\sigma\\
         p\cdot \bar\sigma&m
     \end{pmatrix}\sum_{s=1}^2\xi^s\xi^{\dagger s}.
        \end{aligned}
    \end{equation}
    The latter sum can be obtained by considering
    \begin{equation}
        \sum_{s=1}^2\delta^{sa}\delta^{sb} = \delta^{ab}\implies \sum   _{s=1}^2(\xi ^{\dagger s}\cdot\xi^a)(\xi ^{\dagger b}\cdot\xi^s)  =\xi^{\dagger b}\cdot\xi^a \implies \sum_{s=1}^2\xi^s_\alpha\xi^{\dagger s}_\beta = \delta_{\alpha\beta}.
    \end{equation}
    Therefore, we find
    \begin{equation}
        \sum_{s=1}^{2} u^s(\mathbf{p})\bar{u}^s(\mathbf{p})  = \slashed{p} +m.
    \end{equation}
    Almost identically, we find
    \begin{equation}
        \sum_{s=1}^{2} v^s(\mathbf{p})\bar{v}^s(\mathbf{p}) = \slashed{p} - m.
    \end{equation}
\end{solution}

\subsection{Canonical Quantization of the Dirac Field, Q6}

The Fourier decomposition of the Dirac operator $\psi(\mathbf{x})$ and the conjugate field $\psi^\dagger(\mathbf{x})$ is given by
\begin{align}
    \psi(\mathbf{x}) &= \sum_{s=1}^{2}\int\frac{d^3p}{(2\pi)^3}\frac{1}{\sqrt{2E_{\mathbf{p}}}}\left[b^s_{\mathbf{p}}\,u^s(\mathbf{p})e^{+i\mathbf{p}\cdot\mathbf{x}} + c^{s\dagger}_{\mathbf{p}}\,v^s(\mathbf{p})e^{-i\mathbf{p}\cdot\mathbf{x}}\right]\\
    \psi^\dagger(\mathbf{x}) &= \sum_{s=1}^{2}\int\frac{d^3p}{(2\pi)^3}\frac{1}{\sqrt{2E_{\mathbf{p}}}}\left[b^{s\dagger}_{\mathbf{p}}\,u^s(\mathbf{p})^\dagger e^{-i\mathbf{p}\cdot\mathbf{x}} + c^s_{\mathbf{p}}\,v^s(\mathbf{p})^\dagger e^{+i\mathbf{p}\cdot\mathbf{x}}\right]
\end{align}
The creation and annihilation operators are taken to satisfy
\begin{align}
    \{b^r_{\mathbf{p}}, b^{s\dagger}_{\mathbf{q}}\} &= (2\pi)^3\delta^{rs}\,\delta^{(3)}(\mathbf{p}-\mathbf{q})\\
    \{c^r_{\mathbf{p}}, c^{s\dagger}_{\mathbf{q}}\} &= (2\pi)^3\delta^{rs}\,\delta^{(3)}(\mathbf{p}-\mathbf{q})
\end{align}
with all other anti-commutators vanishing,
\begin{equation}
    \{b^r_{\mathbf{p}}, b^s_{\mathbf{q}}\} = \{c^r_{\mathbf{p}}, c^s_{\mathbf{q}}\} = \{b^r_{\mathbf{p}}, c^{s\dagger}_{\mathbf{q}}\} = \{b^r_{\mathbf{p}}, c^s_{\mathbf{q}}\} = \ldots = 0.
\end{equation}
Show that these imply that the field and its conjugate momenta satisfy the anti-commutation relations,
\begin{align}
    \{\psi_\alpha(\mathbf{x}), \psi_\beta(\mathbf{y})\} &= \{\psi^\dagger_\alpha(\mathbf{x}), \psi^\dagger_\beta(\mathbf{y})\} = 0\\
    \{\psi_\alpha(\mathbf{x}), \psi^\dagger_\beta(\mathbf{y})\} &= \delta_{\alpha\beta}\,\delta^{(3)}(\mathbf{x}-\mathbf{y}).
\end{align}
(Note: The calculation is very similar to that for the bosonic field, but at some point you will need to make use of the identities (\ref{eq:uouter}) and (\ref{eq:vouter})).

\begin{solution}
    We compute $\{\psi_\alpha(\mathbf{x}), \psi^\dagger_\beta(\mathbf{y})\}$ directly. Substituting the Fourier decompositions,
    \begin{equation}
    \begin{aligned}
        \{\psi_\alpha(\mathbf{x}), \psi^\dagger_\beta(\mathbf{y})\} &= \sum_{r,s}\int\frac{d^3p}{(2\pi)^3}\frac{1}{\sqrt{2E_\mathbf{p}}}\int\frac{d^3q}{(2\pi)^3}\frac{1}{\sqrt{2E_\mathbf{q}}}\\
        &\quad\times\bigg[\{b^r_\mathbf{p}, b^{s\dagger}_\mathbf{q}\}u^r_\alpha(\mathbf{p})u^{s\dagger}_\beta(\mathbf{q})e^{i\mathbf{p}\cdot\mathbf{x}-i\mathbf{q}\cdot\mathbf{y}}\\
        &\qquad + \{c^{r\dagger}_\mathbf{p}, c^s_\mathbf{q}\}v^r_\alpha(\mathbf{p})v^{s\dagger}_\beta(\mathbf{q})e^{-i\mathbf{p}\cdot\mathbf{x}+i\mathbf{q}\cdot\mathbf{y}}\bigg]
    \end{aligned}
    \end{equation}
    where all cross terms $\{b,c\}$ and $\{b,c^\dagger\}$ vanish. Then,
    \begin{equation}
    \begin{aligned}
        \{\psi_\alpha(\mathbf{x}), \psi^\dagger_\beta(\mathbf{y})\} &= \int\frac{d^3p}{(2\pi)^3}\frac{1}{2E_\mathbf{p}}\sum_s\left[u^s_\alpha(\mathbf{p})u^{s\dagger}_\beta(\mathbf{p})e^{i\mathbf{p}\cdot(\mathbf{x}-\mathbf{y})} + v^s_\alpha(\mathbf{p})v^{s\dagger}_\beta(\mathbf{p})e^{-i\mathbf{p}\cdot(\mathbf{x}-\mathbf{y})}\right].
    \end{aligned}
    \end{equation}
    Now we use the spin sum identities. Note that $\bar{u} = u^\dagger\gamma^0$ so $u^\dagger = \bar{u}\gamma^0$, giving
    \begin{equation}
        \sum_s u^s u^{s\dagger} = \sum_s u^s\bar{u}^s\gamma^0 = (\slashed{p}+m)\gamma^0, \qquad \sum_s v^s v^{s\dagger} = (\slashed{p}-m)\gamma^0.
    \end{equation}
    Substituting back,
    \begin{equation}
    \begin{aligned}
        \{\psi_\alpha(\mathbf{x}), \psi^\dagger_\beta(\mathbf{y})\} &= \int\frac{d^3p}{(2\pi)^3}\frac{1}{2E_\mathbf{p}}\left[(\slashed{p}+m)_{\alpha\dot\beta}(\gamma^0)_{\dot\beta\beta}e^{i\mathbf{p}\cdot(\mathbf{x}-\mathbf{y})} + (\slashed{p}-m)_{\alpha\dot\beta}(\gamma^0)_{\dot\beta\beta}e^{-i\mathbf{p}\cdot(\mathbf{x}-\mathbf{y})}\right].
    \end{aligned}
    \end{equation}
    Expanding $\slashed{p} = p_0\gamma^0 - p_i\gamma^i$ and using $(\gamma^0)^2 = \mathbf{1}$, the $m$ terms cancel between the two exponentials after relabelling $\mathbf{p}\to-\mathbf{p}$ in the second integral, and the $p_i\gamma^i\gamma^0$ terms also cancel by the same relabelling. What remains is
    \begin{equation}
    \begin{aligned}
        \{\psi_\alpha(\mathbf{x}), \psi^\dagger_\beta(\mathbf{y})\} &= \int\frac{d^3p}{(2\pi)^3}\frac{2p_0}{2E_\mathbf{p}}(\gamma^0\gamma^0)_{\alpha\beta}e^{i\mathbf{p}\cdot(\mathbf{x}-\mathbf{y})}\\
        &= \delta_{\alpha\beta}\int\frac{d^3p}{(2\pi)^3}e^{i\mathbf{p}\cdot(\mathbf{x}-\mathbf{y})} = \delta_{\alpha\beta}\,\delta^{(3)}( {x}-\mathbf{y}).
    \end{aligned}
    \end{equation}
    The vanishing of $\{\psi_\alpha(\mathbf{x}),\psi_\beta(\mathbf{y})\}$ follows similarly, since the cross terms $\{b^r_\mathbf{p}, c^{s\dagger}_\mathbf{q}\} = 0$ and the surviving terms $\{b,b\} = \{c^\dagger, c^\dagger\} = 0$ all vanish.
\end{solution}

\subsection{The Dirac Hamiltonian, Q7}

Using the results of Question 6, show that the quantum Hamiltonian
\begin{equation}
    H = \int d^3x\; \bar\psi(-i\gamma^i\partial_i + m)\psi
\end{equation}
can be written, after normal ordering, as
\begin{equation}
    H = \int\frac{d^3p}{(2\pi)^3}\,E_{\mathbf{p}}\sum_{s=1}^{2}\left[b^{s\dagger}_{\mathbf{p}}b^s_{\mathbf{p}} + c^{s\dagger}_{\mathbf{p}}c^s_{\mathbf{p}}\right].
\end{equation}
(Note: Again, the calculation is very similar to that for the bosonic field. This time you will need to make use of the identities derived in Questions 4 and 5).

\begin{solution}
    Substituting the Fourier decompositions and using $\bar\psi = \psi^\dagger\gamma^0$,
    \begin{equation}
    \begin{aligned}
        H &= \sum_{r,s}\int d^3x\int\frac{d^3p}{(2\pi)^3}\frac{E_\mathbf{p}}{\sqrt{2E_\mathbf{p}}}\int\frac{d^3q}{(2\pi)^3}\frac{1}{\sqrt{2E_\mathbf{q}}}\\
        &\quad\times\left[b^{r\dagger}_\mathbf{q} u^r(\mathbf{q})^\dagger e^{-i\mathbf{q}\cdot\mathbf{x}} + c^r_\mathbf{q} v^r(\mathbf{q})^\dagger e^{+i\mathbf{q}\cdot\mathbf{x}}\right]\left[b^s_\mathbf{p} u^s(\mathbf{p})e^{+i\mathbf{p}\cdot\mathbf{x}} - c^{s\dagger}_\mathbf{p} v^s(\mathbf{p})e^{-i\mathbf{p}\cdot\mathbf{x}}\right]
    \end{aligned}
    \end{equation}
    where we applied $(-i\gamma^i\partial_i + m)u^s(\mathbf{p})e^{i\mathbf{p}\cdot\mathbf{x}} = E_\mathbf{p}\gamma^0 u^s e^{i\mathbf{p}\cdot\mathbf{x}}$ and $(-i\gamma^i\partial_i+m)v^s e^{-i\mathbf{p}\cdot\mathbf{x}} = -E_\mathbf{p}\gamma^0 v^s e^{-i\mathbf{p}\cdot\mathbf{x}}$ from the Dirac equation, and used $(\gamma^0)^2 = \mathbf{1}$ to cancel the $\gamma^0$ from $\bar\psi = \psi^\dagger\gamma^0$. Integrating over $d^3x$ gives delta functions, setting $\mathbf{q}=\mathbf{p}$ for diagonal terms and $\mathbf{q}=-\mathbf{p}$ for cross terms:
    \begin{equation}
    \begin{aligned}
        H &= \sum_{r,s}\int\frac{d^3p}{(2\pi)^3}\frac{E_\mathbf{p}}{2E_\mathbf{p}}\bigg[u^r(\mathbf{p})^\dagger u^s(\mathbf{p})\,b^{r\dagger}_\mathbf{p} b^s_\mathbf{p} - v^r(\mathbf{p})^\dagger v^s(\mathbf{p})\,c^r_\mathbf{p} c^{s\dagger}_\mathbf{p}\\
        &\qquad\qquad\qquad + u^r(-\mathbf{p})^\dagger(-v^s(\mathbf{p}))\,b^{r\dagger}_{-\mathbf{p}}c^{s\dagger}_\mathbf{p} + v^r(-\mathbf{p})^\dagger u^s(\mathbf{p})\,c^r_{-\mathbf{p}}b^s_\mathbf{p}\bigg].
    \end{aligned}
    \end{equation}
    The cross terms vanish by the orthogonality relation $u^r(-\mathbf{p})^\dagger v^s(\mathbf{p}) = 0$. Applying the normalisation $u^r(\mathbf{p})^\dagger u^s(\mathbf{p}) = v^r(\mathbf{p})^\dagger v^s(\mathbf{p}) = 2E_\mathbf{p}\delta^{rs}$ from (5) and (6),
    \begin{equation}
    \begin{aligned}
        H &= \sum_s\int\frac{d^3p}{(2\pi)^3}E_\mathbf{p}\left[b^{s\dagger}_\mathbf{p} b^s_\mathbf{p} - c^s_\mathbf{p} c^{s\dagger}_\mathbf{p}\right]\\
        &= \sum_s\int\frac{d^3p}{(2\pi)^3}E_\mathbf{p}\left[b^{s\dagger}_\mathbf{p} b^s_\mathbf{p} - \left(-c^{s\dagger}_\mathbf{p} c^s_\mathbf{p} + (2\pi)^3\delta^{(3)}(0)\right)\right].
    \end{aligned}
    \end{equation}
    After normal ordering, the vacuum contribution $\delta^{(3)}(0)$ is discarded and the sign of the $c$ term flips, giving
    \begin{equation}
        H = \int\frac{d^3p}{(2\pi)^3}\,E_\mathbf{p}\sum_{s=1}^2\left[b^{s\dagger}_\mathbf{p} b^s_\mathbf{p} + c^{s\dagger}_\mathbf{p} c^s_\mathbf{p}\right].
    \end{equation}
    Note the crucial sign difference from the bosonic case: here normal ordering the anticommuting $c^s_\mathbf{p} c^{s\dagger}_\mathbf{p}$ flips the sign to give a positive contribution, whereas for bosons it produces an infinite additive constant.
\end{solution}

\subsection{The Spin-Statistics Theorem, Q8}

The purpose of this question is to give you a glimpse into the spin-statistics theorem. This theorem roughly says that if you try to quantize a field with the wrong statistics, bad things will happen. Here we'll see what goes wrong if you try to quantize a spin $1/2$ field as a boson. We start with the usual decomposition (11). This time we choose bosonic commutation relations for the annihilation and creation operators,
\begin{align}
    [b^r_{\mathbf{p}}, b^{s\dagger}_{\mathbf{q}}] &= (2\pi)^3\delta^{rs}\,\delta^{(3)}(\mathbf{p}-\mathbf{q})\\
    [c^r_{\mathbf{p}}, c^{s\dagger}_{\mathbf{q}}] &= -(2\pi)^3\delta^{rs}\,\delta^{(3)}(\mathbf{p}-\mathbf{q})
\end{align}
with all other commutators vanishing. Note the strange minus sign for the $c$ operators. Repeat the calculation of Question 6 to show that these are equivalent to the commutation relations,
\begin{align}
    [\psi_\alpha(\mathbf{x}), \psi_\beta(\mathbf{y})] &= [\psi^\dagger_\alpha(\mathbf{x}), \psi^\dagger_\beta(\mathbf{y})] = 0\\
    [\psi_\alpha(\mathbf{x}), \psi^\dagger_\beta(\mathbf{y})] &= \delta_{\alpha\beta}\,\delta^{(3)}(\mathbf{x}-\mathbf{y}).
\end{align}
Now repeat the calculation of Question 7, to show that, after normal ordering, the Hamiltonian is given by
\begin{equation}
    H = \int\frac{d^3p}{(2\pi)^3}\,E_{\mathbf{p}}\sum_{s=1}^{2}\left[b^{s\dagger}_{\mathbf{p}}b^s_{\mathbf{p}} - c^{s\dagger}_{\mathbf{p}}c^s_{\mathbf{p}}\right].
\end{equation}
This Hamiltonian is not bounded below: you can lower the energy indefinitely by creating more and more $c$ particles. This is the reason a theory of bosonic spin $1/2$ particles is sick.

\begin{solution}
    The calculation proceeds exactly as in Questions 6 and 7. The field commutator gives
    \begin{equation}
    \begin{aligned}
        [\psi_\alpha(\mathbf{x}), \psi^\dagger_\beta(\mathbf{y})] &= \int\frac{d^3p}{(2\pi)^3}\frac{1}{2E_\mathbf{p}}\sum_s\left[u^s_\alpha(\mathbf{p})u^{s\dagger}_\beta(\mathbf{p})e^{i\mathbf{p}\cdot(\mathbf{x}-\mathbf{y})} - v^s_\alpha(\mathbf{p})v^{s\dagger}_\beta(\mathbf{p})e^{-i\mathbf{p}\cdot(\mathbf{x}-\mathbf{y})}\right]\\
        &= \int\frac{d^3p}{(2\pi)^3}\frac{1}{2E_\mathbf{p}}\left[(\slashed{p}+m)\gamma^0 e^{i\mathbf{p}\cdot(\mathbf{x}-\mathbf{y})} - (\slashed{p}-m)\gamma^0 e^{-i\mathbf{p}\cdot(\mathbf{x}-\mathbf{y})}\right]_{\alpha\beta}\\
        &= \delta_{\alpha\beta}\,\delta^{(3)}(\mathbf{x}-\mathbf{y})
    \end{aligned}
    \end{equation}
    where the minus sign from $[c,c^\dagger] = -(2\pi)^3\delta^{(3)}$ cancels the minus sign in the $v\bar{v}$ spin sum, reproducing the same result as the fermionic case. For the Hamiltonian,
    \begin{equation}
    \begin{aligned}
        H &= \sum_s\int\frac{d^3p}{(2\pi)^3}E_\mathbf{p}\left[b^{s\dagger}_\mathbf{p} b^s_\mathbf{p} + c^s_\mathbf{p} c^{s\dagger}_\mathbf{p}\right]\\
        &= \sum_s\int\frac{d^3p}{(2\pi)^3}E_\mathbf{p}\left[b^{s\dagger}_\mathbf{p} b^s_\mathbf{p} + c^{s\dagger}_\mathbf{p} c^s_\mathbf{p} + (2\pi)^3\delta^{(3)}(0)\right].
    \end{aligned}
    \end{equation}
    Note that now $c^s_\mathbf{p} c^{s\dagger}_\mathbf{p} = c^{s\dagger}_\mathbf{p} c^s_\mathbf{p} + (2\pi)^3\delta^{(3)}(0)$ since these are bosonic operators. After normal ordering,
    \begin{equation}
        H = \int\frac{d^3p}{(2\pi)^3}\,E_\mathbf{p}\sum_{s=1}^2\left[b^{s\dagger}_\mathbf{p} b^s_\mathbf{p} - c^{s\dagger}_\mathbf{p} c^s_\mathbf{p}\right]
    \end{equation}
    where the minus sign arises from the $[c,c^\dagger]=-(2\pi)^3\delta^{(3)}$ relation flipping the sign under normal ordering.
\end{solution}